% ECONOMICS_PROBLEM: {"id": "E52-395", "title": "Heterogeneous monetary transmission", "jel_code": "E52", "source_id": "OP-0244"}
% !TeX program = xelatex
\documentclass[11pt,a4paper]{article}
\usepackage[margin=23mm]{geometry}
\usepackage{fontspec}
\setmainfont{Latin Modern Roman}
\usepackage{amsmath,amssymb,mathtools}
\usepackage{xcolor,enumitem,booktabs,longtable,array,xurl,hyperref}
\definecolor{muted}{HTML}{53636C}
\hypersetup{colorlinks=true,urlcolor=blue,linkcolor=blue}
\setlength{\parindent}{0pt}
\setlength{\parskip}{5pt}
\newcommand{\R}{\mathbb R}
\newcommand{\E}{\mathbb E}
\newcommand{\N}{\mathbb N}
\newcommand{\ind}{\mathbf 1}
\newcommand{\Var}{\operatorname{Var}}
\newcommand{\Cov}{\operatorname{Cov}}
\newcommand{\supp}{\operatorname{supp}}
\DeclareMathOperator*{\argmin}{arg\,min}
\DeclareMathOperator*{\argmax}{arg\,max}
\newcommand{\entrymeta}[3]{{\sffamily\small\color{muted}\textbf{\texttt{#1}}\quad|\quad #2\par #3\par}}
\newcommand{\Problemref}[1]{\texttt{#1}}
\newcommand{\JELref}[2]{\texttt{#2} (JEL \texttt{#1})}
\begin{document}
\section*{Heterogeneous monetary transmission}
\textbf{Problem ID:} \texttt{E52-395}\quad \textbf{JEL:} \texttt{E52}\par
% BEGIN ECONOMICS STATEMENT
\label{problem:OP-0244}
\label{audited:ECON-025}
{\sffamily\small\color{muted}Original subject group: Inflation and monetary policy\par}
\label{definitions:problem:30007525}
\textbf{Audit classification:} Explicit published open question. The Fisher-channel magnitude is explicitly debated; the conditional causal response surface is editorial. Verification concerns the cited version, not first priority or proof that this exact editorial specification remains unsolved on October 8, 2026.
\subsection*{Definitions and precise research target}
\paragraph{Editorial mathematical specification.} Fix households with predetermined type \(h_i=(a_i^{L},a_i^{I},d_i,n_i)\), where liquid assets, illiquid assets, debt, and nominal asset exposure are measured before a monetary intervention. Let \(m_t\) be an externally identified unexpected increase in the policy rate, and \(c_{i,t+H}(m)\) household real consumption under shock size \(m\), including equilibrium income and price adjustments. Define the conditional response
\[
\tau_H(h)=\left.\frac{\partial E[c_{i,t+H}(m)\mid h_i=h]}{\partial m}\right|_{m=0}.
\]
Identify \(\tau_H(h)\) and its aggregation \(\int\tau_H(h)\,dF(h)\), where \(F\) is the population distribution of predetermined types. Provide instruments that isolate monetary surprises from central-bank information about fundamentals, and address nonrandom refinancing, unobserved insurance, and changes in household composition. Separate direct interest-rate and nominal-revaluation channels from induced labor-income and asset-price responses through a stated structural model or independently credible channel interventions. Test whether identified heterogeneous effects aggregate to observed consumption responses across episodes. The target includes uncertainty about the quantitative Fisher redistribution channel; it does not assume that all transmission mechanisms are unknown or that conditioning on wealth alone identifies a causal channel.

Measure the shock \(m\) in a fixed unit, such as one percentage point of the annualized nominal policy rate, and define consumption in common real currency units. A derivative requires differentiability near zero and two-sided treatment support; without these, target finite shock contrasts instead. For continuous household type \(h\), conditional expectations are defined only \(F\)-almost everywhere. To interchange the derivative and population integral, require an integrable dominating bound on individual derivatives and hold the predetermined type distribution fixed. For a household with nominal net position \(n_i\), a price-level surprise changes its real value locally by \(-n_i\,d\log P\); a consumption channel also needs its marginal propensity to consume and general-equilibrium income response. This balance-sheet identity alone is not a causal estimate.
\subsection*{Variants and assumption changes}
The following are \emph{editorial variants}, not additional published open conjectures.
\begin{enumerate}
\item \emph{Finite shock asymmetry.} Compare \(E[c_H(a)-c_H(0)\mid h]\) and \(E[c_H(-a)-c_H(0)\mid h]\) on overlapping states and treatment supports. Do not infer large-shock symmetry from a local derivative.
\item \emph{Refinancing constraints.} Separate fixed-rate borrowers, adjustable-rate borrowers, and refinancing-eligible borrowers using characteristics measured before the announcement. Compare direct debt-service responses with labor-income responses under a disclosed channel-intervention model.
\end{enumerate}
\subsection*{References and source audit}
\begin{itemize}
\item Janet L. Yellen. \emph{Macroeconomic Research After the Crisis}. 2016. Locator: Section Heterogeneity. Primary text checked. \url{https://www.federalreserve.gov/newsevents/speech/yellen20161014a.htm}.
\item Adrien Auclert; Matthew Rognlie; Ludwig Straub. \emph{Fiscal and Monetary Policy with Heterogeneous Agents}. 2025. Locator: Nominal assets, printed p.16; final sentence on the Fisher-channel empirical magnitude. Primary text checked. \url{https://web.stanford.edu/~aauclert/annual_review_hank.pdf}.
\end{itemize}
The Fisher-channel magnitude is explicitly debated; the conditional causal response surface is editorial.
\begin{enumerate}\raggedright
\item\hypertarget{definitions:source:30007525:1}{} Janet L. Yellen. 2016. \emph{Macroeconomic Research After the Crisis}. Section Heterogeneity.
\url{https://www.federalreserve.gov/newsevents/speech/yellen20161014a.htm}.
\item\hypertarget{definitions:source:30007525:2}{} Adrien Auclert; Matthew Rognlie; Ludwig Straub. 2025. \emph{Fiscal and Monetary Policy with Heterogeneous Agents}. Nominal assets, printed p.16; final sentence on the Fisher-channel empirical magnitude.
\url{https://web.stanford.edu/~aauclert/annual_review_hank.pdf}.
DOI: \url{https://doi.org/10.1146/annurev-economics-091624-044646}.
\end{enumerate}

\subsection*{Integrated refinement: ECON-025}
{\sffamily\small\color{muted}Mortgage structure in heterogeneous monetary transmission\par}
This refinement adopts the additional source conventions
(page~\pageref{audited:conventions}), subject to its local restrictions.
\textbf{Mortgage-contract composition.}
Let the pre-shock state be $s_0=(x_0,\chi)$, with
$x_0=(a,h,b,z)$ and mortgage contract $\chi$, and use the initial law $\mu_0$.
Conditional expected consumption $c_t(s_0;\epsilon)$ integrates subsequent
refinancing, default and state transitions. Define
\[
\operatorname{IRF}_C(t;\mu_0)
=\left.\frac{\partial}{\partial\epsilon}
       \int c_t(s_0;\epsilon)\,d\mu_0(s_0)\right|_{\epsilon=0}.
\]
Compare this response after replacing the supported conditional contract law by
$H'(d\chi\mid x_0)$ while holding the $x_0$ marginal fixed and recomputing
equilibrium. Specify maturity, payment resets, prepayment, collateral and
refinancing eligibility; include lenders and securities holders as counterparts.
Hold shock timing, normalization and the subsequent policy rule fixed. Use finite
impulses when differentiation fails. Reweighting existing contracts is distinct
from changing future menus or ownership selection, and does not freeze the
endogenous date-$t$ state distribution.
\subsection*{Supplied source audit for ECON-025}
Local [S1], [S2], etc. in this audit and the references immediately below
belong to ECON-025, independently of the original entry's reference numbering.
[S1], PDF p. 9, supplies the HANK/housing/mortgage topic. This reweighting estimand is editorial, not an identified formula supplied by the lecture.
\subsection*{References for integrated refinement ECON-025}
{\footnotesize\raggedright
\par\noindent\textbf{[S1]} Benjamin Moll (n.d.). \emph{Lecture 11: Good Research Topics and Open Questions}. Distributional Macroeconomics lecture slides; publicly archived at a 2019 URL.\par \textit{Verified locator / source role:} PDF p. 9: housing, mortgages, and monetary transmission.\par \url{https://benjaminmoll.com/wp-content/uploads/2019/07/Lecture11_2149.pdf}\par\smallskip
}

\subsection*{Common conventions: Source mathematical conventions}

These conventions apply to each entry unless explicitly replaced.

\textbf{Models and identification.} A primitive model \(m\) specifies all probability laws, feasible choices, information, equilibrium and measurement rules used by its target. The admissible class \(\mathcal M\) is fixed independently of outcomes. If \(P_O(m)\) is its observable probability law and \(\tau(m)\) the target, its identified set at observable law \(P\) is
\[
 \mathcal I_\tau(P)=\{\tau(m):m\in\mathcal M,\ P_O(m)=P\}.
\]
Point identification means that this set is a singleton. An empty set signals incompatibility of data and assumptions. Coordinate bounds need not describe the joint set. Bounds are sharp only if every retained target is attainable or arbitrarily approachable by compatible models. Finite bounds require explicit restrictions on unobserved quantities; unrestricted confounding, hidden wealth, extrapolation or arbitrary equilibrium selection can yield uninformative sets.

\textbf{Causal interventions.} A potential outcome \(Y_i(p)\) is the outcome under a complete policy regime \(p\), including timing, financing, information and market spillovers. The target population and units are fixed before treatment. With interference, use the complete assignment vector or a justified exposure mapping; individual no-interference notation alone cannot identify an economy-wide intervention. A difference of mean potential outcomes does not require the cross-world joint distribution. Conditioning on post-treatment migration, survival, enrollment or employment changes the estimand and needs separate justification.

\textbf{Statistical statements.} An estimator, confidence set, forecast or test specifies a sampling law, model class and loss/coverage criterion. Uniform \((1-\alpha)\)-coverage means \(\inf_{m\in\mathcal M}\Pr_m\{\tau(m)\in C_n\}\geq1-\alpha\), or an explicitly asymptotic version. The whole parameter space trivially has coverage, so informativeness requires a stated width, risk or power objective. A forecasting improvement is relative to a named benchmark, information set and evaluation population.

\textbf{Equilibrium and optimization.} An equilibrium comprises feasible plans, optimality, beliefs where needed, budget constraints and all market/resource clearing. The primitives or a stated benchmark define each correspondence \(\mathcal E(p;m)\). Nonemptiness is assumed or proved, never inferred just from compact policy sets. A selected-equilibrium target uses a declared selection rule; a robust target quantifies over all permitted equilibria. Use suprema or infima unless attainment is established. An \(\varepsilon\)-optimal policy is feasible and within \(\varepsilon\) of the finite optimum value. Report an empty feasible set separately.

\textbf{Welfare and accounting.} Normative utility, interpersonal normalization, social weights, numeraire, discounting and the use of public receipts are primitives. Behavioral utility need not coincide with the welfare criterion. Household welfare includes ownership income and rents once; adding profits again to compensating variation generally double-counts them. A monetary equivalent is defined by an explicit transfer and price convention, with monotonicity and range conditions. Average effects, mechanism contributions, transfers and welfare losses are different targets. Infinite sums require convergence and dynamic feasible plans require terminal or no-Ponzi conditions.

\textbf{Source versions.} A working-paper date, revision date and journal publication year are distinguished. A DOI for a journal version is not automatically the DOI of an earlier manuscript. Locators refer to the stated version and printed pages where specified. A metadata-only check does not verify a claim about a full-text conclusion. Every entry's source subsection narrows the provenance claim accordingly.

\subsection*{Common conventions: Additional-source status and mathematical conventions}

\label{audited:conventions}

Of the 100 supplied ECON entries, 65 distinct problems appear here; 32 close
matches refine earlier entries, and three duplicate questions map to existing
entries. Appendix~\ref{app:audited-crosswalk} lists every destination.
The attachment's P/R/S/U distinctions and source audits are retained. Its
precise mathematical formalizations are editorial unless the individual audit
says otherwise. Candidate subsidy targets ECON-004--006 retain their U flags;
the resolved approval-core question ECON-007 updates OP-0202 above.
The following supplied conventions govern both this part and the attributed
ECON refinements elsewhere in the catalogue, subject to local restrictions.

\textbf{Parameterized questions.} A problem family lists inputs that must be fixed before an instance is evaluated: population, horizon, policies, information, data law, model class and welfare weights. These are required choices, not quantities the citations are assumed to specify. Undefined applications should not be treated as identified estimands.

\textbf{Indices and integrability.} Sets of agents, firms and regions are finite unless a population measure is explicitly used. \(i,j\) index units, \(r\) regions, \(t\) dates and \(h\) horizons. \(\mathbb E\) and \(\Pr\) refer to the declared population and law; \(\mathbf1\{A\}\) is an indicator. Every displayed expectation and discounted sum needs existence in the stated domain. Infinite horizons use a discount in \((0,1)\) and a convergence condition; finite horizons may use discount one.

\textbf{Optimization and existence.} A displayed \(\max\), \(\min\), or \(\arg\max\) is conditional on attainment. For finite feasible menus this is immediate; general spaces need continuity and compactness/coercivity or another existence argument. Without attainment, the target is the supremum/infimum and arbitrarily near-optimal feasible choices. \(\inf\varnothing=+\infty\). Empty feasibility and an unidentified target are different issues.

\textbf{Potential outcomes.} \(Y_i(z)\) is the outcome under a specified intervention, not the observed conditional mean among people choosing \(z\). In binary comparisons \(\Delta Y_i=Y_i(1)-Y_i(0)\). Specify assignment, treatment versions, compliance, information and observation horizon. Interference is allowed under a full policy vector; compressed exposure notation needs a justified mapping. No interference, consistency and overlap are substantive assumptions, not automatic consequences of notation.

\textbf{Populations and observation.} Fix baseline eligibility and population weights. Conditioning on post-policy employment, survival, relocation or program uptake can change the population being compared. Death is not ordinary loss to follow-up; outcomes truncated by death need a composite convention, joint survival target, or explicitly defined survivor stratum. Fertility-changing interventions need a population functional rather than an unexplained fixed descendant cohort. Measurement and missingness models must be supplied.

\textbf{Identification.} A structural model \(m\in\mathcal M\) implies observable law \(P_m\) and target \(\theta(m)\). Identification means
\[
 P_{m_1}=P_{m_2}\ \Longrightarrow\ \theta(m_1)=\theta(m_2)
 \quad\text{for all }m_1,m_2\in\mathcal M .
\]
Without identification, the target is
\[
 \Theta(P;\mathcal M)=\{\theta(m):m\in\mathcal M,\ P_m=P\}.
\]
Writing \(\theta(P)\) before establishing identification can hide incompatible latent models. Confidence sets additionally need a sampling law, confidence level, asymptotic sequence and uniformity class. Point identification, coverage and informative precision are separate properties.

\textbf{Mediation.} Total effects of randomized assignment are distinct from cross-world natural indirect effects. Belief/effort-channel effects require a meaningful mediator intervention and additional measurement, exclusion or mediation assumptions. Randomization of the initial policy alone does not supply those assumptions. Report bounds or interventional alternatives when the stronger target is unsupported.

\textbf{Equilibrium.} A counterfactual \(W(z)\), \(p(z)\), or \(\mu_t^z\) requires individual optimization, feasibility, government/resource budgets, market clearing and state-transition laws. State equilibrium existence and selection, or report ranges over compatible equilibria. Initial-state integration includes subsequent endogenous transitions; current-state integration must use the policy-dependent distribution. Reduced-form invariance after policy is not assumed.

\textbf{Welfare accounts.} Scores, utility, health, dollars and ecological quantities require explicit conversion before addition. Fees, taxes, subsidies, wages and extortion payments are transfers between parties; distributional weights can make them welfare relevant, but their face value is not automatically a resource loss. State ownership, geographic boundaries, foreign-income weights and fiscal finance. Do not add profits to household welfare already including dividends, or count land capitalization and the underlying benefit twice. A benefit-cost expression must distinguish gross benefits from net welfare and subtract every real cost exactly once.

\textbf{Derivatives and risk.} Log derivatives require positive arguments; local responses require admissible perturbations and specified equilibrium branches. Differentiation under expectations needs regularity/domination, and boundaries may allow only one-sided derivatives. Risk uses a specified law; ambiguity uses a declared nonempty law set on a common history space. Adaptive policies must be nonanticipative. A robust criterion is an editorial choice unless attributed explicitly.

\textbf{Scope overlaps.} Allocation, collective choice, contracts, household bargaining and coordinated policy overlap economics with optimization and game theory. They are retained because they appeared in the original catalogue; no claim of a strictly game-theory-free selection is made.
% END ECONOMICS STATEMENT
\end{document}
